\section{稳健性、技术批判与取舍}
\subsection{时间异常与距离近似影响什么}
12对时间矛盾中，11对位于inner，1对位于验证，没有落入测试。进一步检查运营时段后，实际训练采样池只少9对，验证的异常对在6:00前，不进入运营池。单独以剔除副本重训并验证，保留原教学baseline；具体收入/应答差异见表\ref{tab:offset}，不能仅因占比小就未经运行宣称完全无影响。
\begin{table}[htbp]\centering\small
\caption{UTC矛盾对的单因素稳健性：60车/300回合、5验证模拟日。}\label{tab:offset}
\begin{tabularx}{\linewidth}{Yrrr}\toprule
训练代理池 & 接单车费 & 应答率 & 接驾/km\\\midrule
教师全代理 & 16,830.68 & 0.5900 & 1.891\\
剔除UTC矛盾对 & 16,669.55 & 0.5849 & 1.905\\
\bottomrule\end{tabularx}
\tabnote{原始25550教学代理保留；剔除副本不改offset、不重排原数据，验证运营池不受那1对6:00前异常影响。}
\end{table}


对25,550个相同坐标对，以WGS84仅作为数学参考的椭球测地线，与教学工作平面的欧氏距离独立比较，后者相对差为$-0.633\%$至$-0.042\%$，最大绝对差约133.50 m。这是固定假设下的近似检查，不能证明源文件datum。更大的差异来自距离模型：曼哈顿/平面欧氏比在1至$\sqrt2$之间，不应把两种度量的差误称投影误差。3 km半径与时间片中的量纲始终保持同一教学模型。

\subsection{独立核验改变了哪些结论依据}
核验从原始实例、OD/批次和逐条派单记录取得参考，没有读取执行者的PASS字符串来代替计算。在线部分以枚举、独立可行性匹配和LP复算两种离线目标；进化部分重解码基因、重加目标、box-union HV和MILP效用参考；RL用原订单流和trace独立回放位置、忙时、取消、派单与收入，并对TD/Double建立手算反例。Q/V是模型估计，不能仅由热图漂亮判断正确。

独立检查真正促成了KM Spacing规范、初始化成本补计、Double随机流分离和图表背景聚合等修正。前两者改变评价有效性或预算解释，随机流修复使受影响Double旧结果失效并重跑；纯绘图背景和版式修正不冒称新实验。这些是不同问题，不能把每一次工程修复画成方法效果上升。

\subsection{为什么停止继续叠加}
多目标的KM初始化在原预算下有明确开发和未来批次支持；额外种群组合只增加少量HV、增加约一倍EA评价并恶化Spacing，没有新的机制证据支持继续扩张。RL中，Double、调度与NSGA融合分别回答了过估计、空间再分布和接驾权衡，但它们没有共同构成全面占优的策略。收入、响应、接驾与空驶出现冲突时，保留各维结果，不在看见结果后发明综合评分。

最终保留的技术认识是：先保证服务覆盖再解释条件延迟；离线参考必须与目标一致；前沿形状与初始化比算法编号更有解释力；长期价值能够改变位置选择，也可能导致过度拒单；移动收益必须与真实占用和成本一起看。结论限定在教学生成机制、签到代理、所列输入和计算预算内，不外推真实出租车平台效果。

\subsection{AI参与与证据限度}
Codex协调真实独立上下文完成来源、执行、批判、制图与交叉核验，研究增强由已批准范围内的实测反馈触发。没有向每条订单调用LLM，也没有把普通优化函数更名为Agent。自动规则内的采用和回退属于Codex执行判断，不能倒写为后来并未发生的人工逐项决定。代码、原始run、负结果、规范指标与报告源保持追溯，外部二审与学生理解仍需另行完成。
